\documentclass[a4paper,11pt,draft]{article}
\usepackage[utf8]{inputenc}
\usepackage[T1]{fontenc}
\usepackage[ngerman]{babel}
\usepackage{amsmath}
\usepackage{graphicx}
\usepackage{geometry}
\usepackage{hyperref}
\usepackage{booktabs}
\usepackage{float}

\geometry{a4paper, left=2.5cm, right=2.5cm, top=2.5cm, bottom=2.5cm}

\title{Protokoll: BLAST-Analyse und Sequenzcharakterisierung von TSR3}
\author{Julian Rasp, Raphael Ortner, Johann Vogel}
\date{\today}

\begin{document}

\maketitle

\section{Einleitung und Proteinidentifikation}
Die durchgeführte BLAST-Analyse (Basic Local Alignment Search Tool) identifizierte die Query-Sequenz eindeutig als \textbf{TSR3} (Ribosome Biogenesis Protein TSR3 Homolog). In den Datenbanken (z.B. NCBI) wird das Protein häufig unter der biochemischen Bezeichnung \textbf{18S rRNA aminocarboxypropyltransferase} (EC

\end{document}